\section{From the thermal inequality to channel outputs}\label{sec:output}

We first turn the thermal matrix inequality into an entropy bound for
$\A_{\gamma,\nu}\otimes\mathcal N$, with any finite-dimensional
partner channel $\mathcal N$. Its two matrices will encode the outputs
of conditional input branches; a common Kraus index will preserve their
cross operator. We then iterate this bound for repeated damping uses,
retaining the input population at every site.

\subsection{The one-qubit entropy and its convexity}

Fix $\gamma,\nu\in[0,1]$ and abbreviate
$\A=\A_{\gamma,\nu}$. Put
\begin{equation}\label{eq:coefficients}
 k=\sqrt{1-\gamma},\qquad b=\gamma(1-\nu),\qquad
 c=\gamma\nu,\qquad a=1-c,\qquad d=1-b.
\end{equation}
Then
\begin{equation}\label{eq:coefficient-identities}
 a+c=b+d=1,\qquad ad-bc=k^2.
\end{equation}
The population transition matrix is
$\bigl(\begin{smallmatrix}a&b\\c&d\end{smallmatrix}\bigr)$ and
the coherence multiplier is $k$. A Kraus representation is
\begin{equation}\label{eq:kraus}
 \begin{aligned}
 L_0&=\sqrt{1-\nu}\,\diag(1,k),&
 L_1&=\sqrt b\,\ket0\bra1,\\
 L_2&=\sqrt\nu\,\diag(k,1),&
 L_3&=\sqrt c\,\ket1\bra0.
 \end{aligned}
\end{equation}
Direct multiplication gives $\A(P)=\sum_{i=0}^3L_iPL_i^*$ and
$\sum_{i=0}^3L_i^*L_i=I$, including all parameter endpoints.

For $p\in[0,1]$, define
\begin{equation}\label{eq:single-entropy}
 \begin{split}
 v(p)&=(a(1-p)+bp)(c(1-p)+dp)-k^2p(1-p)\\
     &=ac(1-p)^2+bdp^2+2bcp(1-p),\qquad e(p)=g(v(p)).
 \end{split}
\end{equation}
Any pure qubit state of excitation probability $p$ has coherence
modulus $\sqrt{p(1-p)}$. Hence $v(p)$ is its output determinant and
$e(p)$ its output entropy. In particular, $v(p)\in[0,1/4]$.
Expanding \eqref{eq:coefficients} in \eqref{eq:single-entropy} gives
\eqref{eq:v-main}. If $\gamma>0$ and $0<\nu<1$,
then $v(p)\ge\gamma\nu(1-\nu)>0$ for every $p$. Every pure input
therefore has a positive definite output. A mixed input is a convex
combination of pure inputs, so its output is positive definite as well.
Thus these thermal parameters lie outside the pure-output class treated
by Tang, Zhu, Bai, and Wang
\cite[Corollary~5.1 and Lemma~5.2]{TangEtAl2026}.

\begin{lemma}\label{lem:e-convex}
For every $\gamma,\nu\in[0,1]$, the function
$p\mapsto g(\gamma\nu(1-\nu)+\gamma(1-\gamma)(p-\nu)^2)$
is convex on $[0,1]$.
\end{lemma}
\begin{proof}
The completed-square formula \eqref{eq:v-main} gives
\[
 \sqrt{v(p)}=
 \left\|\begin{pmatrix}
 \sqrt{\gamma\nu(1-\nu)}\\
 \sqrt{\gamma(1-\gamma)}\,(p-\nu)
 \end{pmatrix}\right\|_2.
\]
This is a convex function of $p$, being the norm of an affine vector,
even when either coefficient vanishes. By
\Cref{lem:binary-spectrum}, $s\mapsto g(s^2)$ is nondecreasing and convex on
$[0,1/2]$, which contains the range of $\sqrt{v}$. A nondecreasing
convex function of a convex function is convex: apply the inner
convexity, then monotonicity, then the outer convexity.
This proves the assertion, with no parameter exclusions.
\end{proof}

\subsection{A single damping use with an arbitrary partner}

The following estimate is the channel form of \Cref{thm:thermal}.
Its scalar term is the pure one-qubit entropy $e(p)$, and the remaining
terms are the output entropies of the two conditional branches.

\begin{lemma}\label{lem:partner-branch}
Let $\gamma,\nu\in[0,1]$, put $\A=\A_{\gamma,\nu}$ and
$e(p)=g(v_{\gamma,\nu}(p))$, and let
$\mathcal N:\mathcal B(\mathcal H_{\mathrm{in}})\to
\mathcal B(\mathcal H_{\mathrm{out}})$ be completely positive and
trace preserving on nonzero finite-dimensional complex Hilbert spaces.
For a unit vector $\psi\in\C^2\otimes\mathcal H_{\mathrm{in}}$, write
\[
 \psi=\sqrt{1-p}\,\ket0\otimes\psi_0
       +\sqrt p\,\ket1\otimes\psi_1,
 \qquad 0\le p\le1,
\]
where each $\psi_i$ is a unit vector and a zero-weight branch may be
chosen arbitrarily. Then
\begin{equation}\label{eq:partner-branch}
 \begin{split}
 \ent\bigl((\A\otimes\mathcal N)(\proj\psi)\bigr)
 &\ge e(p)+(1-p)\ent\bigl(\mathcal N(\proj{\psi_0})\bigr)\\
 &\quad+p\ent\bigl(\mathcal N(\proj{\psi_1})\bigr).
 \end{split}
\end{equation}
\end{lemma}

\begin{proof}
Choose a finite Kraus representation
\[
 \mathcal N(P)=\sum_{\ell=1}^r V_\ell P V_\ell^*,
 \qquad \sum_{\ell=1}^r V_\ell^*V_\ell=I.
\]
Here $V_\ell:\mathcal H_{\mathrm{in}}\to\mathcal H_{\mathrm{out}}$.
Complete positivity in finite dimensions gives such a finite family;
the last identity is trace preservation. If
$d_{\mathrm{out}}=\dim\mathcal H_{\mathrm{out}}$, form
$X,Y\in\C^{d_{\mathrm{out}}\times r}$ by giving their $\ell$th columns
the values $V_\ell\psi_0$ and $V_\ell\psi_1$, respectively. The common
column index gives the three Gram identities and normalizations
\begin{equation}\label{eq:partner-grams}
 \begin{split}
 XX^*&=\mathcal N(\proj{\psi_0}),\qquad
 YY^*=\mathcal N(\proj{\psi_1}),\\
 XY^*&=\mathcal N(\ket{\psi_0}\bra{\psi_1}),\qquad
 \Tr(XX^*)=\Tr(YY^*)=1.
 \end{split}
\end{equation}
No equality between $d_{\mathrm{out}}$, $r$, and the input dimension
is required, and the branch vectors need not be orthogonal.

Set $x=1-p$ and $y=p$, and use $a,b,c,d,k$ from
\eqref{eq:coefficients}. Applying the two channels to the branch
decomposition gives
\[
 (\A\otimes\mathcal N)(\proj\psi)=
 \begin{pmatrix}
 axXX^*+byYY^*&k\sqrt{xy}\,XY^*\\
 k\sqrt{xy}\,YX^*&cxXX^*+dyYY^*
 \end{pmatrix}.
\]
This is \eqref{eq:thermal-state} with
\begin{equation}\label{eq:channel-substitution}
 \alpha=ax,\quad\beta=by,\quad\eta=cx,\quad\delta=dy,
 \qquad K=k\sqrt{xy}.
\end{equation}
The coefficients are nonnegative, sum to one, and satisfy
$\alpha\delta-\beta\eta=(ad-bc)xy=K^2$.
Their column sums are $\alpha+\eta=x$ and $\beta+\delta=y$.
Moreover, their scalar parameter is
\[
 u=(ax+by)(cx+dy)-k^2xy=v(p).
\]
Thus \Cref{thm:thermal} applies to the arbitrary rectangular matrices
$X,Y$ and gives \eqref{eq:partner-branch}. The theorem includes $u=0$,
and the zero-weight branch terms vanish, so this argument includes
$p=0,1$ and every boundary value of $\gamma,\nu$.
\end{proof}

\subsection{Iteration over a block}

We now take the partner to be the channel on the remaining qubits.
Convexity of $e$ will combine the conditional populations at each later
site into that site's original population.

\begin{proposition}\label{prop:pure-bound}
For every $\gamma,\nu\in[0,1]$, put $\A=\A_{\gamma,\nu}$ and
$e(p)=g(v_{\gamma,\nu}(p))$. For every integer $n\ge1$ and
every unit vector $\psi$ on $n$ qubits,
let $p_j$ be the probability of outcome $1$ at site $j$ in the
computational basis. Then
\begin{equation}\label{eq:pure-bound}
 \ent\bigl(\A^{\otimes n}(\proj\psi)\bigr)
 \ge\sum_{j=1}^n e(p_j).
\end{equation}
\end{proposition}

\begin{proof}
For $n=1$ there is equality by \eqref{eq:single-entropy}. Suppose the
claim holds for $n-1$ and write an arbitrary pure input as
\begin{equation}\label{eq:branch-decomposition}
 \psi=\sqrt x\,\ket0\otimes\psi_0
          +\sqrt y\,\ket1\otimes\psi_1,
 \qquad x=1-p_1,\quad y=p_1.
\end{equation}
The branch vectors are normalized but need not be orthogonal. For a
zero-weight branch choose any unit vector. The decomposition is valid
for every input vector, and each conditional vector may be entangled
across the remaining sites.

Apply \Cref{lem:partner-branch} with
$\mathcal N=\A^{\otimes(n-1)}$, and write
$\sigma_i=\A^{\otimes(n-1)}(\proj{\psi_i})$ for the two branch
outputs. The lemma gives the explicit recursion
\begin{equation}\label{eq:recursion}
 \ent\bigl(\A^{\otimes n}(\proj\psi)\bigr)
 \ge e(p_1)+x\ent(\sigma_0)+y\ent(\sigma_1).
\end{equation}
Apply the induction hypothesis to both pure branch inputs. If
$p_{j|i}$ is the excitation probability of site $j>1$ in $\psi_i$,
then
\[
 p_j=xp_{j|0}+yp_{j|1}.
\]
This identity follows from $\langle0|1\rangle=0$ on the first site;
it does not require $\psi_0\perp\psi_1$. By \Cref{lem:e-convex},
\[
 x\ent(\sigma_0)+y\ent(\sigma_1)
 \ge\sum_{j=2}^n\bigl(xe(p_{j|0})+ye(p_{j|1})\bigr)
 \ge\sum_{j=2}^ne(p_j).
\]
Together with \eqref{eq:recursion}, this proves the induction step.
All terms from zero-weight branches have zero coefficient, so the
argument also includes $p_1=0,1$.
\end{proof}

The same population bound holds for mixed inputs. This form will let us
bound every signal in a Holevo ensemble directly.

\begin{corollary}\label{cor:mixed-bound}
For every $\gamma,\nu\in[0,1]$, integer $n\ge1$, and density matrix
$\rho$ on $(\C^2)^{\otimes n}$, let $p_j$ be its excitation
probability at site $j$. Then
\[
 \ent\bigl(\A_{\gamma,\nu}^{\otimes n}(\rho)\bigr)
 \ge\sum_{j=1}^n g(v_{\gamma,\nu}(p_j)).
\]
\end{corollary}
\begin{proof}
Take a finite pure-state decomposition $\rho=\sum_aq_a\proj{\psi_a}$,
and let $p_{a,j}$ be the corresponding populations. Entropy concavity,
\Cref{prop:pure-bound}, and convexity from \Cref{lem:e-convex} give
\[
 \ent\bigl(\A^{\otimes n}(\rho)\bigr)
 \ge\sum_aq_a\sum_j e(p_{a,j})
 \ge\sum_j e\!\left(\sum_aq_ap_{a,j}\right)=\sum_j e(p_j).
\]
\end{proof}
